% from ega0-20-fr.tex
\oldpage[0\textsubscript{IV-1}]{122}
et $C$ sont deux $A$-anneaux topologiques, on notera
$\operatorname{Hom.\ cont}_A(B,C)$ l'ensemble des $A$-homomorphismes continus
de $B$ dans $C$.

Soient $A$ un anneau topologique, $B$ un $A$-anneau topologique, $L$ un
$B$-bimodule topologique ; on note $\operatorname{D\acute{e}r.\ cont}_A(B,L)$
l'ensemble des $A$-dérivations \emph{continues} de $B$ dans $L$ ; il est clair
que c'est un \emph{sous-groupe} de $\operatorname{D\acute{e}r}_A(B,L)$ (et un
sous-$A$-module lorsque $A$ et $B$ sont commutatifs et $L$ un $B$-module).
\end{env}

\begin{env}[20.3.2]
\label{0.20.3.2-fr}
Il est immédiat que la prop. \sref{0.20.1.1-fr} subsiste quand on y remplace
« anneau » par « anneau topologique » et « homomorphisme » par
« homomorphisme continu ». De même, \sref{0.20.1.3-fr} et
\sref{0.20.1.4-fr} restent valables en remplaçant $\operatorname{D\acute{e}r}$ par
$\operatorname{D\acute{e}r.\ cont}$, « anneau » (resp. « bimodule ») par « anneau
topologique » (resp. « bimodule topologique »), « homomorphisme » par
« homomorphisme continu » ; il faut naturellement supposer dans
\sref{0.20.1.4-fr} que $p$ est continu, et remplacer « $A$-équivalences » par
« $A$-équivalences continues ». On a les mêmes résultats pour
\sref{0.20.1.5-fr} et \sref{0.20.1.6-fr}, à condition de prendre comme
topologie sur $D_C(L)$ (resp. $D_B(L)$) la topologie \emph{produit} (sur
$C\times L$, resp. $B\times L$).
\end{env}

\begin{proposition}[20.3.3]
\label{0.20.3.3-fr}
Soient $A$ un anneau topologique, $B$ un $A$-anneau topologique, $L$ un
$B$-bimodule topologique discret annulé par un idéal bilatère ouvert de $B$.
Si dans $B$ le carré de tout idéal bilatère ouvert est ouvert, on a
$\operatorname{D\acute{e}r.\ cont}_A(B,L)=\operatorname{D\acute{e}r}_A(B,L)$.
\end{proposition}

\begin{proof}
En effet, si $\mathfrak{R}$ est un idéal bilatère ouvert de $B$ annulant $L$,
et $D$ une $A$-dérivation de $B$ dans $L$, on a
$D(\mathfrak{R}^2)\subset\mathfrak{R}.D(\mathfrak{R})+
D(\mathfrak{R}).\mathfrak{R}=0$ \sref{0.20.1.1.1-fr}, donc $D$ est continue.
\end{proof}

\begin{env}[20.3.4]
\label{0.20.3.4-fr}
Tous les résultats de \sref{0.20.2.1-fr} restent valables lorsqu'on y remplace
« anneau » par « anneau topologique », « bimodule » par « bimodule
topologique », « homomorphisme » par « homomorphisme continu » et
$\operatorname{D\acute{e}r}$ par $\operatorname{D\acute{e}r.\ cont}$.
\end{env}

\begin{proposition}[20.3.5]
\label{0.20.3.5-fr}
Soient $A$ un anneau topologique, $B$ un $A$-anneau topologique, $L$ un
$B$-bimodule topologique discret annulé par un idéal bilatère ouvert de $B$.
On a alors un isomorphisme canonique
\begin{equation}
\tag{20.3.5.1}
\label{0.20.3.5.1-fr}
  \varinjlim\operatorname{D\acute{e}r}_{A/\mathfrak{J}}
  (B/\mathfrak{R},L)\xrightarrow{\sim}
  \operatorname{D\acute{e}r.\ cont}_A(B,L)
\end{equation}
où dans le premier membre la limite inductive est prise suivant l'ensemble
ordonné filtrant des couples $(\mathfrak{J},\mathfrak{R})$ d'idéaux bilatères
tels que $\mathfrak{R}.L=L.\mathfrak{R}=0$,
$\mathfrak{J}.B\subset\mathfrak{R}$,
$B.\mathfrak{J}\subset\mathfrak{R}$.
\end{proposition}

\begin{proof}
Comme $A/\mathfrak{J}$ et $B/\mathfrak{R}$ sont discrets, on a des
homomorphismes canoniques
$w_{\mathfrak{J},\mathfrak{R}}:
\operatorname{D\acute{e}r}_{A/\mathfrak{J}}(B/\mathfrak{R},L)\to
\operatorname{D\acute{e}r.\ cont}_A(B,L)$ formant un système inductif
\sref{0.20.3.4-fr}, d'où l'homomorphisme \sref{0.20.3.5.1-fr} par passage à
la limite inductive. Comme l'homomorphisme
$B/\mathfrak{R}'\to B/\mathfrak{R}$ est surjectif pour
$\mathfrak{R}\supset\mathfrak{R}'$, il résulte aussitôt de la définition que
l'homomorphisme
$\operatorname{D\acute{e}r}_{A/\mathfrak{J}}(B/\mathfrak{R},L)\to
\operatorname{D\acute{e}r}_{A/\mathfrak{J}}(B/\mathfrak{R}',L)$ (avec
$\mathfrak{R}.L=L.\mathfrak{R}=0$,
$\mathfrak{J}.B\subset\mathfrak{R}'$,
$B.\mathfrak{J}\subset\mathfrak{R}'$) est injectif, et il en est évidemment
de même de l'homomorphisme
$\operatorname{D\acute{e}r}_{A/\mathfrak{J}}(B/\mathfrak{R},L)\to
\operatorname{D\acute{e}r}_{A/\mathfrak{J}'}(B/\mathfrak{R},L)$ pour
$\mathfrak{J}'\subset\mathfrak{J}$ \sref{0.20.2.1-fr} ; on en conclut que
l'homomorphisme \sref{0.20.3.5.1-fr} est injectif. D'autre part, si $D$ est
une $A$-dérivation continue de $B$ dans $L$, son noyau contient un idéal
bilatère ouvert $\mathfrak{R}_0$ de $B$, et si $\mathfrak{J}_0$ est un idéal
bilatère ouvert de $A$ tel que $\mathfrak{J}_0B\subset\mathfrak{R}_0$ et
$B\mathfrak{J}_0\subset\mathfrak{R}_0$, il est clair que $D$ est l'image
canonique d'une $(A/\mathfrak{J}_0)$-dérivation de
$B/\mathfrak{R}_0$ dans $L$, donc \sref{0.20.3.5.1-fr} est surjectif.
\end{proof}

